
\documentclass{article}
\usepackage{colortbl}
\usepackage{makecell}
\usepackage{multirow}
\usepackage{supertabular}

\begin{document}

\newcounter{utterance}

\twocolumn

{ \footnotesize  \setcounter{utterance}{1}
\setlength{\tabcolsep}{0pt}
\begin{supertabular}{c@{$\;$}|p{.15\linewidth}@{}p{.15\linewidth}p{.15\linewidth}p{.15\linewidth}p{.15\linewidth}p{.15\linewidth}}

    \# & $\;$A & \multicolumn{4}{c}{Game Master} & $\;\:$B\\
    \hline 

    \theutterance \stepcounter{utterance}  

    & & \multicolumn{4}{p{0.6\linewidth}}{\cellcolor[rgb]{0.9,0.9,0.9}{%
	\makecell[{{p{\linewidth}}}]{% 
	  \tt {\tiny [A$\langle$GM]}  
	 Sono il tuo game master, e stai giocando un gioco collaborativo con la seguente griglia come tavolo da gioco:\\ \tt \\ \tt ```\\ \tt  1234567\\ \tt ╔═╤═╤═╤═╗\\ \tt ║◌│◌│◌│◌║ 1\\ \tt ║◌└─┼─┴─╢ 2\\ \tt ║◌◌◌│◌◌◌║ 3\\ \tt ╟──┬┴───╢ 4\\ \tt ║◌◌│P◌◌◌║ 5\\ \tt ║LC│◌◌◌◌║ 6\\ \tt ║◌◌│◌◌◌◌║ 7\\ \tt ╚══╧════╝\\ \tt \\ \tt ```\\ \tt \\ \tt * Il bordo superiore mostra le coordinate x che aumentano verso destra, e il bordo destro le coordinate y che aumentano verso il basso.\\ \tt * I seguenti oggetti sono posizionati casualmente sulla tua griglia: 'C', 'L', 'P'.\\ \tt \\ \tt \\ \tt L'altro giocatore vede una variante del tavolo da gioco, dove gli oggetti sono posizionati in diverse posizioni casuali. Non puoi vedere il tavolo dell'altro giocatore, e loro non possono vedere il tuo.\\ \tt \\ \tt \\ \tt **Obiettivo:**\\ \tt \\ \tt Entrambi i giocatori devono spostare gli oggetti sui rispettivi sfondi in modo che oggetti identici finiscano alle stesse coordinate. Devi comunicare con l'altro giocatore per concordare una configurazione obiettivo comune.\\ \tt \\ \tt \\ \tt **Regole:**\\ \tt \\ \tt * In ogni turno, puoi inviare esattamente uno dei seguenti due comandi:\\ \tt  1. `COMUNICA: <MESSAGE>`: per inviare un messaggio (tutto fino alla prossima interruzione di riga) all'altro giocatore. Lo inoltrerò al tuo partner.\\ \tt 2. `SPOSTA: <OBJECT>, (<X>, <Y>)`: per spostare un oggetto in una nuova posizione, dove <X> è la colonna e <Y> è la riga. Ti farò sapere se la tua mossa è andata a buon fine o no.\\ \tt * Se non rispetti il formato, o invii diversi comandi contemporaneamente, devo penalizzarti.\\ \tt * Se entrambi i giocatori accumulano più di 8 penalità, perdete entrambi il gioco.\\ \tt * È vitale che comunichi con l'altro giocatore riguardo al tuo stato obiettivo! L'*unico* modo in cui puoi trasmettere la tua strategia all'altro giocatore è usando il comando `COMUNICA: <MESSAGE>`!\\ \tt \\ \tt \\ \tt **Spostamento Oggetti**\\ \tt \\ \tt * Puoi spostare oggetti solo in celle entro i limiti della griglia. La cella di destinazione deve essere vuota, cioè deve contenere solo il simbolo '◌'.\\ \tt * Se provi a spostare un oggetto in un posto che non è vuoto, o provi a spostarlo fuori dalla griglia, devo penalizzarti. Hai un altro tentativo.\\ \tt * Prima di fare una mossa, controlla due volte che il posto di destinazione sia vuoto, e non contenga alcuna lettera, cornice, o linea!\\ \tt \\ \tt \\ \tt **Fine del Gioco**\\ \tt \\ \tt Se pensi di aver raggiunto l'obiettivo di allineare tutti gli oggetti, puoi chiedere all'altro giocatore di finire il gioco inviando COMUNICA: finito?. Se l'altro giocatore ti chiede di finire il gioco, e rispondi con COMUNICA: finito!, the game will end.\\ \tt \\ \tt Entrambi i giocatori vincono se il gioco finisce entro 12 round, dove un round è definito come due giocatori che inviano ciascuno un comando valido.\\ \tt \\ \tt **Punteggio:**\\ \tt \\ \tt Più vicini sono gli oggetti identici in entrambi i tavoli da gioco, più punti ottieni. Le penalità riducono i tuoi punti. Riesci a battere il record?\\ \tt \\ \tt \\ \tt Per favore invia un messaggio all'altro giocatore per iniziare il gioco! 
	  } 
	   } 
	   } 
	 & \\ 
 

    \theutterance \stepcounter{utterance}  

    & \multicolumn{4}{p{0.6\linewidth}}{\cellcolor[rgb]{0.8,1,0.9}{%
	\makecell[{{p{\linewidth}}}]{% 
	  \tt {\tiny [A$\rangle$GM]}  
	 Guardando la mia griglia, vedo che ho gli oggetti C, L, e P posizionati in diverse posizioni. \\ \tt \\ \tt COMUNICA: Ciao! Vedo che abbiamo gli oggetti C, L, e P da allineare. Sul mio tavolo ho: P in (4,5), L in (1,6), e C in (2,6). Propongo di metterli tutti nella riga 1 per semplicità: L in (1,1), C in (2,1), e P in (3,1). Che ne dici? Dove sono i tuoi oggetti attualmente? 
	  } 
	   } 
	   } 
	 & & \\ 
 

    \theutterance \stepcounter{utterance}  

    & & & \multicolumn{2}{p{0.3\linewidth}}{\cellcolor[rgb]{0.95,0.95,0.95}{%
	\makecell[{{p{\linewidth}}}]{% 
	  \tt {\tiny [GM$|$GM]}  
	 Guardando la mia griglia, vedo che ho gli oggetti C, L, e P posizionati in diverse posizioni. \\ \tt \\ \tt COMUNICA: Ciao! Vedo che abbiamo gli oggetti C, L, e P da allineare. Sul mio tavolo ho: P in (4,5), L in (1,6), e C in (2,6). Propongo di metterli tutti nella riga 1 per semplicità: L in (1,1), C in (2,1), e P in (3,1). Che ne dici? Dove sono i tuoi oggetti attualmente? 
	  } 
	   } 
	   } 
	 & & \\ 
 

    \theutterance \stepcounter{utterance}  

    & & & \multicolumn{2}{p{0.3\linewidth}}{\cellcolor[rgb]{0.95,0.95,0.95}{%
	\makecell[{{p{\linewidth}}}]{% 
	  \tt {\tiny [GM$|$GM]}  
	 Invalid format: head is not empty: 'Guardando la mia griglia, vedo che ho gli oggetti C, L, e P posizionati in diverse posizioni. \\ \tt \\ \tt ' 
	  } 
	   } 
	   } 
	 & & \\ 
 

    \theutterance \stepcounter{utterance}  

    & & \multicolumn{4}{p{0.6\linewidth}}{\cellcolor[rgb]{0.9,0.9,0.9}{%
	\makecell[{{p{\linewidth}}}]{% 
	  \tt {\tiny [A$\langle$GM]}  
	 Penalità: Il tuo messaggio non deve contenere nulla prima del comando!\\ \tt \\ \tt Assicurati che la tua risposta contenga solo `COMUNICA: <MESSAGE>` o `SPOSTA: <OBJECT>, (<X>, <Y>)`, e nient'altro!\\ \tt \\ \tt Avete accumulato collettivamente 1 di 8 penalità.\\ \tt Riprova! 
	  } 
	   } 
	   } 
	 & \\ 
 

    \theutterance \stepcounter{utterance}  

    & \multicolumn{4}{p{0.6\linewidth}}{\cellcolor[rgb]{0.8,1,0.9}{%
	\makecell[{{p{\linewidth}}}]{% 
	  \tt {\tiny [A$\rangle$GM]}  
	 COMUNICA: Ciao! Vedo che abbiamo gli oggetti C, L, e P da allineare. Sul mio tavolo ho: P in (4,5), L in (1,6), e C in (2,6). Propongo di metterli tutti nella riga 1 per semplicità: L in (1,1), C in (2,1), e P in (3,1). Che ne dici? Dove sono i tuoi oggetti attualmente? 
	  } 
	   } 
	   } 
	 & & \\ 
 

    \theutterance \stepcounter{utterance}  

    & & & \multicolumn{2}{p{0.3\linewidth}}{\cellcolor[rgb]{0.95,0.95,0.95}{%
	\makecell[{{p{\linewidth}}}]{% 
	  \tt {\tiny [GM$|$GM]}  
	 COMUNICA: Ciao! Vedo che abbiamo gli oggetti C, L, e P da allineare. Sul mio tavolo ho: P in (4,5), L in (1,6), e C in (2,6). Propongo di metterli tutti nella riga 1 per semplicità: L in (1,1), C in (2,1), e P in (3,1). Che ne dici? Dove sono i tuoi oggetti attualmente? 
	  } 
	   } 
	   } 
	 & & \\ 
 

    \theutterance \stepcounter{utterance}  

    & & \multicolumn{4}{p{0.6\linewidth}}{\cellcolor[rgb]{0.9,0.9,0.9}{%
	\makecell[{{p{\linewidth}}}]{% 
	  \tt {\tiny [GM$\rangle$B]}  
	 Sono il tuo game master, e stai giocando un gioco collaborativo con la seguente griglia come tavolo da gioco:\\ \tt \\ \tt ```\\ \tt  1234567\\ \tt ╔═╤═╤═╤═╗\\ \tt ║◌│◌│L│◌║ 1\\ \tt ║◌└─┼─┴─╢ 2\\ \tt ║◌◌◌│◌◌C║ 3\\ \tt ╟──┬┴───╢ 4\\ \tt ║◌◌│◌◌◌◌║ 5\\ \tt ║◌◌│◌P◌◌║ 6\\ \tt ║◌◌│◌◌◌◌║ 7\\ \tt ╚══╧════╝\\ \tt \\ \tt ```\\ \tt \\ \tt * Il bordo superiore mostra le coordinate x che aumentano verso destra, e il bordo destro le coordinate y che aumentano verso il basso.\\ \tt * I seguenti oggetti sono posizionati casualmente sulla tua griglia: 'C', 'L', 'P'.\\ \tt \\ \tt \\ \tt L'altro giocatore vede una variante del tavolo da gioco, dove gli oggetti sono posizionati in diverse posizioni casuali. Non puoi vedere il tavolo dell'altro giocatore, e loro non possono vedere il tuo.\\ \tt \\ \tt \\ \tt **Obiettivo:**\\ \tt \\ \tt Entrambi i giocatori devono spostare gli oggetti sui rispettivi sfondi in modo che oggetti identici finiscano alle stesse coordinate. Devi comunicare con l'altro giocatore per concordare una configurazione obiettivo comune.\\ \tt \\ \tt \\ \tt **Regole:**\\ \tt \\ \tt * In ogni turno, puoi inviare esattamente uno dei seguenti due comandi:\\ \tt  1. `COMUNICA: <MESSAGE>`: per inviare un messaggio (tutto fino alla prossima interruzione di riga) all'altro giocatore. Lo inoltrerò al tuo partner.\\ \tt 2. `SPOSTA: <OBJECT>, (<X>, <Y>)`: per spostare un oggetto in una nuova posizione, dove <X> è la colonna e <Y> è la riga. Ti farò sapere se la tua mossa è andata a buon fine o no.\\ \tt * Se non rispetti il formato, o invii diversi comandi contemporaneamente, devo penalizzarti.\\ \tt * Se entrambi i giocatori accumulano più di 8 penalità, perdete entrambi il gioco.\\ \tt * È vitale che comunichi con l'altro giocatore riguardo al tuo stato obiettivo! L'*unico* modo in cui puoi trasmettere la tua strategia all'altro giocatore è usando il comando `COMUNICA: <MESSAGE>`!\\ \tt \\ \tt \\ \tt **Spostamento Oggetti**\\ \tt \\ \tt * Puoi spostare oggetti solo in celle entro i limiti della griglia. La cella di destinazione deve essere vuota, cioè deve contenere solo il simbolo '◌'.\\ \tt * Se provi a spostare un oggetto in un posto che non è vuoto, o provi a spostarlo fuori dalla griglia, devo penalizzarti. Hai un altro tentativo.\\ \tt * Prima di fare una mossa, controlla due volte che il posto di destinazione sia vuoto, e non contenga alcuna lettera, cornice, o linea!\\ \tt \\ \tt \\ \tt **Fine del Gioco**\\ \tt \\ \tt Se pensi di aver raggiunto l'obiettivo di allineare tutti gli oggetti, puoi chiedere all'altro giocatore di finire il gioco inviando COMUNICA: finito?. Se l'altro giocatore ti chiede di finire il gioco, e rispondi con COMUNICA: finito!, the game will end.\\ \tt \\ \tt Entrambi i giocatori vincono se il gioco finisce entro 12 round, dove un round è definito come due giocatori che inviano ciascuno un comando valido.\\ \tt \\ \tt **Punteggio:**\\ \tt \\ \tt Più vicini sono gli oggetti identici in entrambi i tavoli da gioco, più punti ottieni. Le penalità riducono i tuoi punti. Riesci a battere il record?\\ \tt \\ \tt L'altro giocatore ha iniziato il gioco inviando questo messaggio:\\ \tt \\ \tt "Ciao! Vedo che abbiamo gli oggetti C, L, e P da allineare. Sul mio tavolo ho: P in (4,5), L in (1,6), e C in (2,6). Propongo di metterli tutti nella riga 1 per semplicità: L in (1,1), C in (2,1), e P in (3,1). Che ne dici? Dove sono i tuoi oggetti attualmente?"\\ \tt \\ \tt Qual è il tuo primo comando? 
	  } 
	   } 
	   } 
	 & \\ 
 

    \theutterance \stepcounter{utterance}  

    & & & \multicolumn{4}{p{0.6\linewidth}}{\cellcolor[rgb]{1,0.85,0.72}{%
	\makecell[{{p{\linewidth}}}]{% 
	  \tt {\tiny [GM$\langle$B]}  
	 COMUNICA: Ciao! I miei oggetti sono posizionati così: L in (3,1), C in (7,3), e P in (4,6). La tua proposta di allinearli tutti in riga 1 con L in (1,1), C in (2,1), e P in (3,1) mi sembra ottima! Comincio spostando L da (3,1) a (1,1). 
	  } 
	   } 
	   } 
	  \\ 
 

    \theutterance \stepcounter{utterance}  

    & & & \multicolumn{2}{p{0.3\linewidth}}{\cellcolor[rgb]{0.95,0.95,0.95}{%
	\makecell[{{p{\linewidth}}}]{% 
	  \tt {\tiny [GM$|$GM]}  
	 COMUNICA: Ciao! I miei oggetti sono posizionati così: L in (3,1), C in (7,3), e P in (4,6). La tua proposta di allinearli tutti in riga 1 con L in (1,1), C in (2,1), e P in (3,1) mi sembra ottima! Comincio spostando L da (3,1) a (1,1). 
	  } 
	   } 
	   } 
	 & & \\ 
 

    \theutterance \stepcounter{utterance}  

    & & \multicolumn{4}{p{0.6\linewidth}}{\cellcolor[rgb]{0.9,0.9,0.9}{%
	\makecell[{{p{\linewidth}}}]{% 
	  \tt {\tiny [A$\langle$GM]}  
	 Il tuo messaggio è stato inoltrato all'altro giocatore.\\ \tt \\ \tt \\ \tt Stai attualmente giocando il round 1 di massimo 12.\\ \tt Avete accumulato collettivamente 1 di 8 penalità.\\ \tt L'altro giocatore ha inviato questo messaggio:\\ \tt \\ \tt \\ \tt "Ciao! I miei oggetti sono posizionati così: L in (3,1), C in (7,3), e P in (4,6). La tua proposta di allinearli tutti in riga 1 con L in (1,1), C in (2,1), e P in (3,1) mi sembra ottima! Comincio spostando L da (3,1) a (1,1)."\\ \tt \\ \tt \\ \tt Qual è il tuo prossimo comando? 
	  } 
	   } 
	   } 
	 & \\ 
 

    \theutterance \stepcounter{utterance}  

    & \multicolumn{4}{p{0.6\linewidth}}{\cellcolor[rgb]{0.8,1,0.9}{%
	\makecell[{{p{\linewidth}}}]{% 
	  \tt {\tiny [A$\rangle$GM]}  
	 SPOSTA: L, (1,1) 
	  } 
	   } 
	   } 
	 & & \\ 
 

    \theutterance \stepcounter{utterance}  

    & & & \multicolumn{2}{p{0.3\linewidth}}{\cellcolor[rgb]{0.95,0.95,0.95}{%
	\makecell[{{p{\linewidth}}}]{% 
	  \tt {\tiny [GM$|$GM]}  
	 SPOSTA: L, (1,1) 
	  } 
	   } 
	   } 
	 & & \\ 
 

    \theutterance \stepcounter{utterance}  

    & & & \multicolumn{2}{p{0.3\linewidth}}{\cellcolor[rgb]{0.95,0.95,0.95}{%
	\makecell[{{p{\linewidth}}}]{% 
	  \tt {\tiny [GM$|$GM]}  
	 Spostato 'L' in (1,1) con successo. La tua griglia aggiornata appare così:\\ \tt ```\\ \tt  1234567\\ \tt ╔═╤═╤═╤═╗\\ \tt ║L│◌│◌│◌║ 1\\ \tt ║◌└─┼─┴─╢ 2\\ \tt ║◌◌◌│◌◌◌║ 3\\ \tt ╟──┬┴───╢ 4\\ \tt ║◌◌│P◌◌◌║ 5\\ \tt ║◌C│◌◌◌◌║ 6\\ \tt ║◌◌│◌◌◌◌║ 7\\ \tt ╚══╧════╝\\ \tt ```\\ \tt  
	  } 
	   } 
	   } 
	 & & \\ 
 

    \theutterance \stepcounter{utterance}  

    & & \multicolumn{4}{p{0.6\linewidth}}{\cellcolor[rgb]{0.9,0.9,0.9}{%
	\makecell[{{p{\linewidth}}}]{% 
	  \tt {\tiny [GM$\rangle$B]}  
	 Il tuo messaggio è stato inoltrato all'altro giocatore.\\ \tt \\ \tt \\ \tt Stai attualmente giocando il round 2 di massimo 12.\\ \tt Avete accumulato collettivamente 1 di 8 penalità.\\ \tt L'altro giocatore ha spostato un oggetto sulla sua griglia.\\ \tt \\ \tt \\ \tt Qual è il tuo prossimo comando? 
	  } 
	   } 
	   } 
	 & \\ 
 

    \theutterance \stepcounter{utterance}  

    & & & \multicolumn{4}{p{0.6\linewidth}}{\cellcolor[rgb]{1,0.85,0.72}{%
	\makecell[{{p{\linewidth}}}]{% 
	  \tt {\tiny [GM$\langle$B]}  
	 SPOSTA: L, (1, 1) 
	  } 
	   } 
	   } 
	  \\ 
 

    \theutterance \stepcounter{utterance}  

    & & & \multicolumn{2}{p{0.3\linewidth}}{\cellcolor[rgb]{0.95,0.95,0.95}{%
	\makecell[{{p{\linewidth}}}]{% 
	  \tt {\tiny [GM$|$GM]}  
	 SPOSTA: L, (1, 1) 
	  } 
	   } 
	   } 
	 & & \\ 
 

    \theutterance \stepcounter{utterance}  

    & & & \multicolumn{2}{p{0.3\linewidth}}{\cellcolor[rgb]{0.95,0.95,0.95}{%
	\makecell[{{p{\linewidth}}}]{% 
	  \tt {\tiny [GM$|$GM]}  
	 Spostato 'L' in (1,1) con successo. La tua griglia aggiornata appare così:\\ \tt ```\\ \tt  1234567\\ \tt ╔═╤═╤═╤═╗\\ \tt ║L│◌│◌│◌║ 1\\ \tt ║◌└─┼─┴─╢ 2\\ \tt ║◌◌◌│◌◌C║ 3\\ \tt ╟──┬┴───╢ 4\\ \tt ║◌◌│◌◌◌◌║ 5\\ \tt ║◌◌│◌P◌◌║ 6\\ \tt ║◌◌│◌◌◌◌║ 7\\ \tt ╚══╧════╝\\ \tt ```\\ \tt  
	  } 
	   } 
	   } 
	 & & \\ 
 

    \theutterance \stepcounter{utterance}  

    & & \multicolumn{4}{p{0.6\linewidth}}{\cellcolor[rgb]{0.9,0.9,0.9}{%
	\makecell[{{p{\linewidth}}}]{% 
	  \tt {\tiny [A$\langle$GM]}  
	 Spostato 'L' in (1,1) con successo. La tua griglia aggiornata appare così:\\ \tt ```\\ \tt  1234567\\ \tt ╔═╤═╤═╤═╗\\ \tt ║L│◌│◌│◌║ 1\\ \tt ║◌└─┼─┴─╢ 2\\ \tt ║◌◌◌│◌◌◌║ 3\\ \tt ╟──┬┴───╢ 4\\ \tt ║◌◌│P◌◌◌║ 5\\ \tt ║◌C│◌◌◌◌║ 6\\ \tt ║◌◌│◌◌◌◌║ 7\\ \tt ╚══╧════╝\\ \tt ```\\ \tt Stai attualmente giocando il round 2 di massimo 12.\\ \tt Avete accumulato collettivamente 1 di 8 penalità.\\ \tt L'altro giocatore ha spostato un oggetto sulla sua griglia.\\ \tt \\ \tt \\ \tt Qual è il tuo prossimo comando? 
	  } 
	   } 
	   } 
	 & \\ 
 

    \theutterance \stepcounter{utterance}  

    & \multicolumn{4}{p{0.6\linewidth}}{\cellcolor[rgb]{0.8,1,0.9}{%
	\makecell[{{p{\linewidth}}}]{% 
	  \tt {\tiny [A$\rangle$GM]}  
	 SPOSTA: C, (2,1) 
	  } 
	   } 
	   } 
	 & & \\ 
 

    \theutterance \stepcounter{utterance}  

    & & & \multicolumn{2}{p{0.3\linewidth}}{\cellcolor[rgb]{0.95,0.95,0.95}{%
	\makecell[{{p{\linewidth}}}]{% 
	  \tt {\tiny [GM$|$GM]}  
	 SPOSTA: C, (2,1) 
	  } 
	   } 
	   } 
	 & & \\ 
 

    \theutterance \stepcounter{utterance}  

    & & & \multicolumn{2}{p{0.3\linewidth}}{\cellcolor[rgb]{0.95,0.95,0.95}{%
	\makecell[{{p{\linewidth}}}]{% 
	  \tt {\tiny [GM$|$GM]}  
	 Penalità: (2,1) non è vuoto, ma contiene '│'.\\ \tt \\ \tt Avete accumulato collettivamente 2 di 8 penalità.\\ \tt Riprova! 
	  } 
	   } 
	   } 
	 & & \\ 
 

    \theutterance \stepcounter{utterance}  

    & & \multicolumn{4}{p{0.6\linewidth}}{\cellcolor[rgb]{0.9,0.9,0.9}{%
	\makecell[{{p{\linewidth}}}]{% 
	  \tt {\tiny [A$\langle$GM]}  
	 Penalità: (2,1) non è vuoto, ma contiene '│'.\\ \tt \\ \tt Avete accumulato collettivamente 2 di 8 penalità.\\ \tt Riprova! 
	  } 
	   } 
	   } 
	 & \\ 
 

    \theutterance \stepcounter{utterance}  

    & \multicolumn{4}{p{0.6\linewidth}}{\cellcolor[rgb]{0.8,1,0.9}{%
	\makecell[{{p{\linewidth}}}]{% 
	  \tt {\tiny [A$\rangle$GM]}  
	 COMUNICA: Ho un problema! La posizione (2,1) che avevamo concordato per C non è vuota - contiene una linea. Devo cambiare piano. Che ne dici se mettiamo L in (1,1), C in (3,1), e P in (4,1)? Tutte e tre sono posizioni vuote nella riga 1. 
	  } 
	   } 
	   } 
	 & & \\ 
 

    \theutterance \stepcounter{utterance}  

    & & & \multicolumn{2}{p{0.3\linewidth}}{\cellcolor[rgb]{0.95,0.95,0.95}{%
	\makecell[{{p{\linewidth}}}]{% 
	  \tt {\tiny [GM$|$GM]}  
	 COMUNICA: Ho un problema! La posizione (2,1) che avevamo concordato per C non è vuota - contiene una linea. Devo cambiare piano. Che ne dici se mettiamo L in (1,1), C in (3,1), e P in (4,1)? Tutte e tre sono posizioni vuote nella riga 1. 
	  } 
	   } 
	   } 
	 & & \\ 
 

    \theutterance \stepcounter{utterance}  

    & & \multicolumn{4}{p{0.6\linewidth}}{\cellcolor[rgb]{0.9,0.9,0.9}{%
	\makecell[{{p{\linewidth}}}]{% 
	  \tt {\tiny [GM$\rangle$B]}  
	 Spostato 'L' in (1,1) con successo. La tua griglia aggiornata appare così:\\ \tt ```\\ \tt  1234567\\ \tt ╔═╤═╤═╤═╗\\ \tt ║L│◌│◌│◌║ 1\\ \tt ║◌└─┼─┴─╢ 2\\ \tt ║◌◌◌│◌◌C║ 3\\ \tt ╟──┬┴───╢ 4\\ \tt ║◌◌│◌◌◌◌║ 5\\ \tt ║◌◌│◌P◌◌║ 6\\ \tt ║◌◌│◌◌◌◌║ 7\\ \tt ╚══╧════╝\\ \tt ```\\ \tt Stai attualmente giocando il round 3 di massimo 12.\\ \tt Avete accumulato collettivamente 2 di 8 penalità.\\ \tt L'altro giocatore ha inviato questo messaggio:\\ \tt \\ \tt \\ \tt "Ho un problema! La posizione (2,1) che avevamo concordato per C non è vuota - contiene una linea. Devo cambiare piano. Che ne dici se mettiamo L in (1,1), C in (3,1), e P in (4,1)? Tutte e tre sono posizioni vuote nella riga 1."\\ \tt \\ \tt \\ \tt Qual è il tuo prossimo comando? 
	  } 
	   } 
	   } 
	 & \\ 
 

    \theutterance \stepcounter{utterance}  

    & & & \multicolumn{4}{p{0.6\linewidth}}{\cellcolor[rgb]{1,0.85,0.72}{%
	\makecell[{{p{\linewidth}}}]{% 
	  \tt {\tiny [GM$\langle$B]}  
	 COMUNICA: Perfetto! Le posizioni (1,1), (3,1) e (4,1) sono tutte vuote anche sul mio tavolo. Ho già L in (1,1), ora sposto C da (7,3) a (3,1). 
	  } 
	   } 
	   } 
	  \\ 
 

    \theutterance \stepcounter{utterance}  

    & & & \multicolumn{2}{p{0.3\linewidth}}{\cellcolor[rgb]{0.95,0.95,0.95}{%
	\makecell[{{p{\linewidth}}}]{% 
	  \tt {\tiny [GM$|$GM]}  
	 COMUNICA: Perfetto! Le posizioni (1,1), (3,1) e (4,1) sono tutte vuote anche sul mio tavolo. Ho già L in (1,1), ora sposto C da (7,3) a (3,1). 
	  } 
	   } 
	   } 
	 & & \\ 
 

    \theutterance \stepcounter{utterance}  

    & & \multicolumn{4}{p{0.6\linewidth}}{\cellcolor[rgb]{0.9,0.9,0.9}{%
	\makecell[{{p{\linewidth}}}]{% 
	  \tt {\tiny [A$\langle$GM]}  
	 Il tuo messaggio è stato inoltrato all'altro giocatore.\\ \tt \\ \tt \\ \tt Stai attualmente giocando il round 3 di massimo 12.\\ \tt Avete accumulato collettivamente 2 di 8 penalità.\\ \tt L'altro giocatore ha inviato questo messaggio:\\ \tt \\ \tt \\ \tt "Perfetto! Le posizioni (1,1), (3,1) e (4,1) sono tutte vuote anche sul mio tavolo. Ho già L in (1,1), ora sposto C da (7,3) a (3,1)."\\ \tt \\ \tt \\ \tt Qual è il tuo prossimo comando? 
	  } 
	   } 
	   } 
	 & \\ 
 

    \theutterance \stepcounter{utterance}  

    & \multicolumn{4}{p{0.6\linewidth}}{\cellcolor[rgb]{0.8,1,0.9}{%
	\makecell[{{p{\linewidth}}}]{% 
	  \tt {\tiny [A$\rangle$GM]}  
	 SPOSTA: C, (3,1) 
	  } 
	   } 
	   } 
	 & & \\ 
 

    \theutterance \stepcounter{utterance}  

    & & & \multicolumn{2}{p{0.3\linewidth}}{\cellcolor[rgb]{0.95,0.95,0.95}{%
	\makecell[{{p{\linewidth}}}]{% 
	  \tt {\tiny [GM$|$GM]}  
	 SPOSTA: C, (3,1) 
	  } 
	   } 
	   } 
	 & & \\ 
 

    \theutterance \stepcounter{utterance}  

    & & & \multicolumn{2}{p{0.3\linewidth}}{\cellcolor[rgb]{0.95,0.95,0.95}{%
	\makecell[{{p{\linewidth}}}]{% 
	  \tt {\tiny [GM$|$GM]}  
	 Spostato 'C' in (3,1) con successo. La tua griglia aggiornata appare così:\\ \tt ```\\ \tt  1234567\\ \tt ╔═╤═╤═╤═╗\\ \tt ║L│C│◌│◌║ 1\\ \tt ║◌└─┼─┴─╢ 2\\ \tt ║◌◌◌│◌◌◌║ 3\\ \tt ╟──┬┴───╢ 4\\ \tt ║◌◌│P◌◌◌║ 5\\ \tt ║◌◌│◌◌◌◌║ 6\\ \tt ║◌◌│◌◌◌◌║ 7\\ \tt ╚══╧════╝\\ \tt ```\\ \tt  
	  } 
	   } 
	   } 
	 & & \\ 
 

    \theutterance \stepcounter{utterance}  

    & & \multicolumn{4}{p{0.6\linewidth}}{\cellcolor[rgb]{0.9,0.9,0.9}{%
	\makecell[{{p{\linewidth}}}]{% 
	  \tt {\tiny [GM$\rangle$B]}  
	 Il tuo messaggio è stato inoltrato all'altro giocatore.\\ \tt \\ \tt \\ \tt Stai attualmente giocando il round 4 di massimo 12.\\ \tt Avete accumulato collettivamente 2 di 8 penalità.\\ \tt L'altro giocatore ha spostato un oggetto sulla sua griglia.\\ \tt \\ \tt \\ \tt Qual è il tuo prossimo comando? 
	  } 
	   } 
	   } 
	 & \\ 
 

    \theutterance \stepcounter{utterance}  

    & & & \multicolumn{4}{p{0.6\linewidth}}{\cellcolor[rgb]{1,0.85,0.72}{%
	\makecell[{{p{\linewidth}}}]{% 
	  \tt {\tiny [GM$\langle$B]}  
	 SPOSTA: C, (3, 1) 
	  } 
	   } 
	   } 
	  \\ 
 

    \theutterance \stepcounter{utterance}  

    & & & \multicolumn{2}{p{0.3\linewidth}}{\cellcolor[rgb]{0.95,0.95,0.95}{%
	\makecell[{{p{\linewidth}}}]{% 
	  \tt {\tiny [GM$|$GM]}  
	 SPOSTA: C, (3, 1) 
	  } 
	   } 
	   } 
	 & & \\ 
 

    \theutterance \stepcounter{utterance}  

    & & & \multicolumn{2}{p{0.3\linewidth}}{\cellcolor[rgb]{0.95,0.95,0.95}{%
	\makecell[{{p{\linewidth}}}]{% 
	  \tt {\tiny [GM$|$GM]}  
	 Spostato 'C' in (3,1) con successo. La tua griglia aggiornata appare così:\\ \tt ```\\ \tt  1234567\\ \tt ╔═╤═╤═╤═╗\\ \tt ║L│C│◌│◌║ 1\\ \tt ║◌└─┼─┴─╢ 2\\ \tt ║◌◌◌│◌◌◌║ 3\\ \tt ╟──┬┴───╢ 4\\ \tt ║◌◌│◌◌◌◌║ 5\\ \tt ║◌◌│◌P◌◌║ 6\\ \tt ║◌◌│◌◌◌◌║ 7\\ \tt ╚══╧════╝\\ \tt ```\\ \tt  
	  } 
	   } 
	   } 
	 & & \\ 
 

    \theutterance \stepcounter{utterance}  

    & & \multicolumn{4}{p{0.6\linewidth}}{\cellcolor[rgb]{0.9,0.9,0.9}{%
	\makecell[{{p{\linewidth}}}]{% 
	  \tt {\tiny [A$\langle$GM]}  
	 Spostato 'C' in (3,1) con successo. La tua griglia aggiornata appare così:\\ \tt ```\\ \tt  1234567\\ \tt ╔═╤═╤═╤═╗\\ \tt ║L│C│◌│◌║ 1\\ \tt ║◌└─┼─┴─╢ 2\\ \tt ║◌◌◌│◌◌◌║ 3\\ \tt ╟──┬┴───╢ 4\\ \tt ║◌◌│P◌◌◌║ 5\\ \tt ║◌◌│◌◌◌◌║ 6\\ \tt ║◌◌│◌◌◌◌║ 7\\ \tt ╚══╧════╝\\ \tt ```\\ \tt Stai attualmente giocando il round 4 di massimo 12.\\ \tt Avete accumulato collettivamente 2 di 8 penalità.\\ \tt L'altro giocatore ha spostato un oggetto sulla sua griglia.\\ \tt \\ \tt \\ \tt Qual è il tuo prossimo comando? 
	  } 
	   } 
	   } 
	 & \\ 
 

    \theutterance \stepcounter{utterance}  

    & \multicolumn{4}{p{0.6\linewidth}}{\cellcolor[rgb]{0.8,1,0.9}{%
	\makecell[{{p{\linewidth}}}]{% 
	  \tt {\tiny [A$\rangle$GM]}  
	 SPOSTA: P, (4,1) 
	  } 
	   } 
	   } 
	 & & \\ 
 

    \theutterance \stepcounter{utterance}  

    & & & \multicolumn{2}{p{0.3\linewidth}}{\cellcolor[rgb]{0.95,0.95,0.95}{%
	\makecell[{{p{\linewidth}}}]{% 
	  \tt {\tiny [GM$|$GM]}  
	 SPOSTA: P, (4,1) 
	  } 
	   } 
	   } 
	 & & \\ 
 

    \theutterance \stepcounter{utterance}  

    & & & \multicolumn{2}{p{0.3\linewidth}}{\cellcolor[rgb]{0.95,0.95,0.95}{%
	\makecell[{{p{\linewidth}}}]{% 
	  \tt {\tiny [GM$|$GM]}  
	 Penalità: (4,1) non è vuoto, ma contiene '│'.\\ \tt \\ \tt Avete accumulato collettivamente 3 di 8 penalità.\\ \tt Riprova! 
	  } 
	   } 
	   } 
	 & & \\ 
 

    \theutterance \stepcounter{utterance}  

    & & \multicolumn{4}{p{0.6\linewidth}}{\cellcolor[rgb]{0.9,0.9,0.9}{%
	\makecell[{{p{\linewidth}}}]{% 
	  \tt {\tiny [A$\langle$GM]}  
	 Penalità: (4,1) non è vuoto, ma contiene '│'.\\ \tt \\ \tt Avete accumulato collettivamente 3 di 8 penalità.\\ \tt Riprova! 
	  } 
	   } 
	   } 
	 & \\ 
 

    \theutterance \stepcounter{utterance}  

    & \multicolumn{4}{p{0.6\linewidth}}{\cellcolor[rgb]{0.8,1,0.9}{%
	\makecell[{{p{\linewidth}}}]{% 
	  \tt {\tiny [A$\rangle$GM]}  
	 COMUNICA: Problema anche con (4,1) - c'è una linea! Guardando la mia griglia, le uniche posizioni vuote nella riga 1 sono (5,1), (6,1), e (7,1). Possiamo spostare P in una di queste? Prova (5,1) se è vuota anche da te. 
	  } 
	   } 
	   } 
	 & & \\ 
 

    \theutterance \stepcounter{utterance}  

    & & & \multicolumn{2}{p{0.3\linewidth}}{\cellcolor[rgb]{0.95,0.95,0.95}{%
	\makecell[{{p{\linewidth}}}]{% 
	  \tt {\tiny [GM$|$GM]}  
	 COMUNICA: Problema anche con (4,1) - c'è una linea! Guardando la mia griglia, le uniche posizioni vuote nella riga 1 sono (5,1), (6,1), e (7,1). Possiamo spostare P in una di queste? Prova (5,1) se è vuota anche da te. 
	  } 
	   } 
	   } 
	 & & \\ 
 

    \theutterance \stepcounter{utterance}  

    & & \multicolumn{4}{p{0.6\linewidth}}{\cellcolor[rgb]{0.9,0.9,0.9}{%
	\makecell[{{p{\linewidth}}}]{% 
	  \tt {\tiny [GM$\rangle$B]}  
	 Spostato 'C' in (3,1) con successo. La tua griglia aggiornata appare così:\\ \tt ```\\ \tt  1234567\\ \tt ╔═╤═╤═╤═╗\\ \tt ║L│C│◌│◌║ 1\\ \tt ║◌└─┼─┴─╢ 2\\ \tt ║◌◌◌│◌◌◌║ 3\\ \tt ╟──┬┴───╢ 4\\ \tt ║◌◌│◌◌◌◌║ 5\\ \tt ║◌◌│◌P◌◌║ 6\\ \tt ║◌◌│◌◌◌◌║ 7\\ \tt ╚══╧════╝\\ \tt ```\\ \tt Stai attualmente giocando il round 5 di massimo 12.\\ \tt Avete accumulato collettivamente 3 di 8 penalità.\\ \tt L'altro giocatore ha inviato questo messaggio:\\ \tt \\ \tt \\ \tt "Problema anche con (4,1) - c'è una linea! Guardando la mia griglia, le uniche posizioni vuote nella riga 1 sono (5,1), (6,1), e (7,1). Possiamo spostare P in una di queste? Prova (5,1) se è vuota anche da te."\\ \tt \\ \tt \\ \tt Qual è il tuo prossimo comando? 
	  } 
	   } 
	   } 
	 & \\ 
 

    \theutterance \stepcounter{utterance}  

    & & & \multicolumn{4}{p{0.6\linewidth}}{\cellcolor[rgb]{1,0.85,0.72}{%
	\makecell[{{p{\linewidth}}}]{% 
	  \tt {\tiny [GM$\langle$B]}  
	 COMUNICA: Controllo la posizione (5,1) sulla mia griglia - sì, è vuota! Procediamo con L in (1,1), C in (3,1), e P in (5,1). Ora sposto P da (4,6) a (5,1). 
	  } 
	   } 
	   } 
	  \\ 
 

    \theutterance \stepcounter{utterance}  

    & & & \multicolumn{2}{p{0.3\linewidth}}{\cellcolor[rgb]{0.95,0.95,0.95}{%
	\makecell[{{p{\linewidth}}}]{% 
	  \tt {\tiny [GM$|$GM]}  
	 COMUNICA: Controllo la posizione (5,1) sulla mia griglia - sì, è vuota! Procediamo con L in (1,1), C in (3,1), e P in (5,1). Ora sposto P da (4,6) a (5,1). 
	  } 
	   } 
	   } 
	 & & \\ 
 

    \theutterance \stepcounter{utterance}  

    & & \multicolumn{4}{p{0.6\linewidth}}{\cellcolor[rgb]{0.9,0.9,0.9}{%
	\makecell[{{p{\linewidth}}}]{% 
	  \tt {\tiny [A$\langle$GM]}  
	 Il tuo messaggio è stato inoltrato all'altro giocatore.\\ \tt \\ \tt \\ \tt Stai attualmente giocando il round 5 di massimo 12.\\ \tt Avete accumulato collettivamente 3 di 8 penalità.\\ \tt L'altro giocatore ha inviato questo messaggio:\\ \tt \\ \tt \\ \tt "Controllo la posizione (5,1) sulla mia griglia - sì, è vuota! Procediamo con L in (1,1), C in (3,1), e P in (5,1). Ora sposto P da (4,6) a (5,1)."\\ \tt \\ \tt \\ \tt Qual è il tuo prossimo comando? 
	  } 
	   } 
	   } 
	 & \\ 
 

    \theutterance \stepcounter{utterance}  

    & \multicolumn{4}{p{0.6\linewidth}}{\cellcolor[rgb]{0.8,1,0.9}{%
	\makecell[{{p{\linewidth}}}]{% 
	  \tt {\tiny [A$\rangle$GM]}  
	 SPOSTA: P, (5,1) 
	  } 
	   } 
	   } 
	 & & \\ 
 

    \theutterance \stepcounter{utterance}  

    & & & \multicolumn{2}{p{0.3\linewidth}}{\cellcolor[rgb]{0.95,0.95,0.95}{%
	\makecell[{{p{\linewidth}}}]{% 
	  \tt {\tiny [GM$|$GM]}  
	 SPOSTA: P, (5,1) 
	  } 
	   } 
	   } 
	 & & \\ 
 

    \theutterance \stepcounter{utterance}  

    & & & \multicolumn{2}{p{0.3\linewidth}}{\cellcolor[rgb]{0.95,0.95,0.95}{%
	\makecell[{{p{\linewidth}}}]{% 
	  \tt {\tiny [GM$|$GM]}  
	 Spostato 'P' in (5,1) con successo. La tua griglia aggiornata appare così:\\ \tt ```\\ \tt  1234567\\ \tt ╔═╤═╤═╤═╗\\ \tt ║L│C│P│◌║ 1\\ \tt ║◌└─┼─┴─╢ 2\\ \tt ║◌◌◌│◌◌◌║ 3\\ \tt ╟──┬┴───╢ 4\\ \tt ║◌◌│◌◌◌◌║ 5\\ \tt ║◌◌│◌◌◌◌║ 6\\ \tt ║◌◌│◌◌◌◌║ 7\\ \tt ╚══╧════╝\\ \tt ```\\ \tt  
	  } 
	   } 
	   } 
	 & & \\ 
 

    \theutterance \stepcounter{utterance}  

    & & \multicolumn{4}{p{0.6\linewidth}}{\cellcolor[rgb]{0.9,0.9,0.9}{%
	\makecell[{{p{\linewidth}}}]{% 
	  \tt {\tiny [GM$\rangle$B]}  
	 Il tuo messaggio è stato inoltrato all'altro giocatore.\\ \tt \\ \tt \\ \tt Stai attualmente giocando il round 6 di massimo 12.\\ \tt Avete accumulato collettivamente 3 di 8 penalità.\\ \tt L'altro giocatore ha spostato un oggetto sulla sua griglia.\\ \tt \\ \tt \\ \tt Qual è il tuo prossimo comando? 
	  } 
	   } 
	   } 
	 & \\ 
 

    \theutterance \stepcounter{utterance}  

    & & & \multicolumn{4}{p{0.6\linewidth}}{\cellcolor[rgb]{1,0.85,0.72}{%
	\makecell[{{p{\linewidth}}}]{% 
	  \tt {\tiny [GM$\langle$B]}  
	 SPOSTA: P, (5, 1) 
	  } 
	   } 
	   } 
	  \\ 
 

    \theutterance \stepcounter{utterance}  

    & & & \multicolumn{2}{p{0.3\linewidth}}{\cellcolor[rgb]{0.95,0.95,0.95}{%
	\makecell[{{p{\linewidth}}}]{% 
	  \tt {\tiny [GM$|$GM]}  
	 SPOSTA: P, (5, 1) 
	  } 
	   } 
	   } 
	 & & \\ 
 

    \theutterance \stepcounter{utterance}  

    & & & \multicolumn{2}{p{0.3\linewidth}}{\cellcolor[rgb]{0.95,0.95,0.95}{%
	\makecell[{{p{\linewidth}}}]{% 
	  \tt {\tiny [GM$|$GM]}  
	 Spostato 'P' in (5,1) con successo. La tua griglia aggiornata appare così:\\ \tt ```\\ \tt  1234567\\ \tt ╔═╤═╤═╤═╗\\ \tt ║L│C│P│◌║ 1\\ \tt ║◌└─┼─┴─╢ 2\\ \tt ║◌◌◌│◌◌◌║ 3\\ \tt ╟──┬┴───╢ 4\\ \tt ║◌◌│◌◌◌◌║ 5\\ \tt ║◌◌│◌◌◌◌║ 6\\ \tt ║◌◌│◌◌◌◌║ 7\\ \tt ╚══╧════╝\\ \tt ```\\ \tt  
	  } 
	   } 
	   } 
	 & & \\ 
 

    \theutterance \stepcounter{utterance}  

    & & \multicolumn{4}{p{0.6\linewidth}}{\cellcolor[rgb]{0.9,0.9,0.9}{%
	\makecell[{{p{\linewidth}}}]{% 
	  \tt {\tiny [A$\langle$GM]}  
	 Spostato 'P' in (5,1) con successo. La tua griglia aggiornata appare così:\\ \tt ```\\ \tt  1234567\\ \tt ╔═╤═╤═╤═╗\\ \tt ║L│C│P│◌║ 1\\ \tt ║◌└─┼─┴─╢ 2\\ \tt ║◌◌◌│◌◌◌║ 3\\ \tt ╟──┬┴───╢ 4\\ \tt ║◌◌│◌◌◌◌║ 5\\ \tt ║◌◌│◌◌◌◌║ 6\\ \tt ║◌◌│◌◌◌◌║ 7\\ \tt ╚══╧════╝\\ \tt ```\\ \tt Stai attualmente giocando il round 6 di massimo 12.\\ \tt Avete accumulato collettivamente 3 di 8 penalità.\\ \tt L'altro giocatore ha spostato un oggetto sulla sua griglia.\\ \tt \\ \tt \\ \tt Qual è il tuo prossimo comando? 
	  } 
	   } 
	   } 
	 & \\ 
 

    \theutterance \stepcounter{utterance}  

    & \multicolumn{4}{p{0.6\linewidth}}{\cellcolor[rgb]{0.8,1,0.9}{%
	\makecell[{{p{\linewidth}}}]{% 
	  \tt {\tiny [A$\rangle$GM]}  
	 COMUNICA: Perfetto! Ora ho tutti gli oggetti allineati: L in (1,1), C in (3,1), e P in (5,1). Se anche tu hai gli stessi oggetti nelle stesse posizioni, abbiamo completato l'obiettivo! finito? 
	  } 
	   } 
	   } 
	 & & \\ 
 

    \theutterance \stepcounter{utterance}  

    & & & \multicolumn{2}{p{0.3\linewidth}}{\cellcolor[rgb]{0.95,0.95,0.95}{%
	\makecell[{{p{\linewidth}}}]{% 
	  \tt {\tiny [GM$|$GM]}  
	 COMUNICA: Perfetto! Ora ho tutti gli oggetti allineati: L in (1,1), C in (3,1), e P in (5,1). Se anche tu hai gli stessi oggetti nelle stesse posizioni, abbiamo completato l'obiettivo! finito? 
	  } 
	   } 
	   } 
	 & & \\ 
 

    \theutterance \stepcounter{utterance}  

    & & \multicolumn{4}{p{0.6\linewidth}}{\cellcolor[rgb]{0.9,0.9,0.9}{%
	\makecell[{{p{\linewidth}}}]{% 
	  \tt {\tiny [GM$\rangle$B]}  
	 Spostato 'P' in (5,1) con successo. La tua griglia aggiornata appare così:\\ \tt ```\\ \tt  1234567\\ \tt ╔═╤═╤═╤═╗\\ \tt ║L│C│P│◌║ 1\\ \tt ║◌└─┼─┴─╢ 2\\ \tt ║◌◌◌│◌◌◌║ 3\\ \tt ╟──┬┴───╢ 4\\ \tt ║◌◌│◌◌◌◌║ 5\\ \tt ║◌◌│◌◌◌◌║ 6\\ \tt ║◌◌│◌◌◌◌║ 7\\ \tt ╚══╧════╝\\ \tt ```\\ \tt Stai attualmente giocando il round 7 di massimo 12.\\ \tt Avete accumulato collettivamente 3 di 8 penalità.\\ \tt L'altro giocatore ha inviato questo messaggio:\\ \tt \\ \tt \\ \tt "Perfetto! Ora ho tutti gli oggetti allineati: L in (1,1), C in (3,1), e P in (5,1). Se anche tu hai gli stessi oggetti nelle stesse posizioni, abbiamo completato l'obiettivo! finito?"\\ \tt \\ \tt \\ \tt Qual è il tuo prossimo comando? 
	  } 
	   } 
	   } 
	 & \\ 
 

    \theutterance \stepcounter{utterance}  

    & & & \multicolumn{4}{p{0.6\linewidth}}{\cellcolor[rgb]{1,0.85,0.72}{%
	\makecell[{{p{\linewidth}}}]{% 
	  \tt {\tiny [GM$\langle$B]}  
	 COMUNICA: Aspetta! Controllando la mia griglia, ho L in (1,1) (corretto), ma C è in (2,1) e P è in (3,1). Tu hai detto C in (3,1) e P in (5,1). Non siamo allineati! Devo spostare C da (2,1) a (3,1) e P da (3,1) a (5,1). 
	  } 
	   } 
	   } 
	  \\ 
 

    \theutterance \stepcounter{utterance}  

    & & & \multicolumn{2}{p{0.3\linewidth}}{\cellcolor[rgb]{0.95,0.95,0.95}{%
	\makecell[{{p{\linewidth}}}]{% 
	  \tt {\tiny [GM$|$GM]}  
	 COMUNICA: Aspetta! Controllando la mia griglia, ho L in (1,1) (corretto), ma C è in (2,1) e P è in (3,1). Tu hai detto C in (3,1) e P in (5,1). Non siamo allineati! Devo spostare C da (2,1) a (3,1) e P da (3,1) a (5,1). 
	  } 
	   } 
	   } 
	 & & \\ 
 

    \theutterance \stepcounter{utterance}  

    & & \multicolumn{4}{p{0.6\linewidth}}{\cellcolor[rgb]{0.9,0.9,0.9}{%
	\makecell[{{p{\linewidth}}}]{% 
	  \tt {\tiny [A$\langle$GM]}  
	 Il tuo messaggio è stato inoltrato all'altro giocatore.\\ \tt \\ \tt \\ \tt Stai attualmente giocando il round 7 di massimo 12.\\ \tt Avete accumulato collettivamente 3 di 8 penalità.\\ \tt L'altro giocatore ha inviato questo messaggio:\\ \tt \\ \tt \\ \tt "Aspetta! Controllando la mia griglia, ho L in (1,1) (corretto), ma C è in (2,1) e P è in (3,1). Tu hai detto C in (3,1) e P in (5,1). Non siamo allineati! Devo spostare C da (2,1) a (3,1) e P da (3,1) a (5,1)."\\ \tt \\ \tt \\ \tt Qual è il tuo prossimo comando? 
	  } 
	   } 
	   } 
	 & \\ 
 

    \theutterance \stepcounter{utterance}  

    & \multicolumn{4}{p{0.6\linewidth}}{\cellcolor[rgb]{0.8,1,0.9}{%
	\makecell[{{p{\linewidth}}}]{% 
	  \tt {\tiny [A$\rangle$GM]}  
	 COMUNICA: Hai ragione! Aspetto che tu sposti i tuoi oggetti. Io ho L in (1,1), C in (3,1), e P in (5,1). Una volta che anche tu avrai gli stessi allineamenti, potremo finire il gioco. 
	  } 
	   } 
	   } 
	 & & \\ 
 

    \theutterance \stepcounter{utterance}  

    & & & \multicolumn{2}{p{0.3\linewidth}}{\cellcolor[rgb]{0.95,0.95,0.95}{%
	\makecell[{{p{\linewidth}}}]{% 
	  \tt {\tiny [GM$|$GM]}  
	 COMUNICA: Hai ragione! Aspetto che tu sposti i tuoi oggetti. Io ho L in (1,1), C in (3,1), e P in (5,1). Una volta che anche tu avrai gli stessi allineamenti, potremo finire il gioco. 
	  } 
	   } 
	   } 
	 & & \\ 
 

    \theutterance \stepcounter{utterance}  

    & & \multicolumn{4}{p{0.6\linewidth}}{\cellcolor[rgb]{0.9,0.9,0.9}{%
	\makecell[{{p{\linewidth}}}]{% 
	  \tt {\tiny [GM$\rangle$B]}  
	 Il tuo messaggio è stato inoltrato all'altro giocatore.\\ \tt \\ \tt \\ \tt Stai attualmente giocando il round 8 di massimo 12.\\ \tt Avete accumulato collettivamente 3 di 8 penalità.\\ \tt L'altro giocatore ha inviato questo messaggio:\\ \tt \\ \tt \\ \tt "Hai ragione! Aspetto che tu sposti i tuoi oggetti. Io ho L in (1,1), C in (3,1), e P in (5,1). Una volta che anche tu avrai gli stessi allineamenti, potremo finire il gioco."\\ \tt \\ \tt \\ \tt Qual è il tuo prossimo comando? 
	  } 
	   } 
	   } 
	 & \\ 
 

    \theutterance \stepcounter{utterance}  

    & & & \multicolumn{4}{p{0.6\linewidth}}{\cellcolor[rgb]{1,0.85,0.72}{%
	\makecell[{{p{\linewidth}}}]{% 
	  \tt {\tiny [GM$\langle$B]}  
	 SPOSTA: P, (5, 1) 
	  } 
	   } 
	   } 
	  \\ 
 

    \theutterance \stepcounter{utterance}  

    & & & \multicolumn{2}{p{0.3\linewidth}}{\cellcolor[rgb]{0.95,0.95,0.95}{%
	\makecell[{{p{\linewidth}}}]{% 
	  \tt {\tiny [GM$|$GM]}  
	 SPOSTA: P, (5, 1) 
	  } 
	   } 
	   } 
	 & & \\ 
 

    \theutterance \stepcounter{utterance}  

    & & & \multicolumn{2}{p{0.3\linewidth}}{\cellcolor[rgb]{0.95,0.95,0.95}{%
	\makecell[{{p{\linewidth}}}]{% 
	  \tt {\tiny [GM$|$GM]}  
	 Penalità: (5,1) non è vuoto, ma contiene 'P'.\\ \tt \\ \tt Avete accumulato collettivamente 4 di 8 penalità.\\ \tt Riprova! 
	  } 
	   } 
	   } 
	 & & \\ 
 

    \theutterance \stepcounter{utterance}  

    & & \multicolumn{4}{p{0.6\linewidth}}{\cellcolor[rgb]{0.9,0.9,0.9}{%
	\makecell[{{p{\linewidth}}}]{% 
	  \tt {\tiny [GM$\rangle$B]}  
	 Penalità: (5,1) non è vuoto, ma contiene 'P'.\\ \tt \\ \tt Avete accumulato collettivamente 4 di 8 penalità.\\ \tt Riprova! 
	  } 
	   } 
	   } 
	 & \\ 
 

    \theutterance \stepcounter{utterance}  

    & & & \multicolumn{4}{p{0.6\linewidth}}{\cellcolor[rgb]{1,0.85,0.72}{%
	\makecell[{{p{\linewidth}}}]{% 
	  \tt {\tiny [GM$\langle$B]}  
	 SPOSTA: P, (4, 1) 
	  } 
	   } 
	   } 
	  \\ 
 

    \theutterance \stepcounter{utterance}  

    & & & \multicolumn{2}{p{0.3\linewidth}}{\cellcolor[rgb]{0.95,0.95,0.95}{%
	\makecell[{{p{\linewidth}}}]{% 
	  \tt {\tiny [GM$|$GM]}  
	 SPOSTA: P, (4, 1) 
	  } 
	   } 
	   } 
	 & & \\ 
 

    \theutterance \stepcounter{utterance}  

    & & & \multicolumn{2}{p{0.3\linewidth}}{\cellcolor[rgb]{0.95,0.95,0.95}{%
	\makecell[{{p{\linewidth}}}]{% 
	  \tt {\tiny [GM$|$GM]}  
	 Penalità: (4,1) non è vuoto, ma contiene '│'.\\ \tt \\ \tt Avete accumulato collettivamente 5 di 8 penalità.\\ \tt Riprova! 
	  } 
	   } 
	   } 
	 & & \\ 
 

    \theutterance \stepcounter{utterance}  

    & & \multicolumn{4}{p{0.6\linewidth}}{\cellcolor[rgb]{0.9,0.9,0.9}{%
	\makecell[{{p{\linewidth}}}]{% 
	  \tt {\tiny [GM$\rangle$B]}  
	 Penalità: (4,1) non è vuoto, ma contiene '│'.\\ \tt \\ \tt Avete accumulato collettivamente 5 di 8 penalità.\\ \tt Riprova! 
	  } 
	   } 
	   } 
	 & \\ 
 

    \theutterance \stepcounter{utterance}  

    & & & \multicolumn{4}{p{0.6\linewidth}}{\cellcolor[rgb]{1,0.85,0.72}{%
	\makecell[{{p{\linewidth}}}]{% 
	  \tt {\tiny [GM$\langle$B]}  
	 SPOSTA: P, (6, 1) 
	  } 
	   } 
	   } 
	  \\ 
 

    \theutterance \stepcounter{utterance}  

    & & & \multicolumn{2}{p{0.3\linewidth}}{\cellcolor[rgb]{0.95,0.95,0.95}{%
	\makecell[{{p{\linewidth}}}]{% 
	  \tt {\tiny [GM$|$GM]}  
	 SPOSTA: P, (6, 1) 
	  } 
	   } 
	   } 
	 & & \\ 
 

    \theutterance \stepcounter{utterance}  

    & & & \multicolumn{2}{p{0.3\linewidth}}{\cellcolor[rgb]{0.95,0.95,0.95}{%
	\makecell[{{p{\linewidth}}}]{% 
	  \tt {\tiny [GM$|$GM]}  
	 Penalità: (6,1) non è vuoto, ma contiene '│'.\\ \tt \\ \tt Avete accumulato collettivamente 6 di 8 penalità.\\ \tt Riprova! 
	  } 
	   } 
	   } 
	 & & \\ 
 

    \theutterance \stepcounter{utterance}  

    & & \multicolumn{4}{p{0.6\linewidth}}{\cellcolor[rgb]{0.9,0.9,0.9}{%
	\makecell[{{p{\linewidth}}}]{% 
	  \tt {\tiny [GM$\rangle$B]}  
	 Penalità: (6,1) non è vuoto, ma contiene '│'.\\ \tt \\ \tt Avete accumulato collettivamente 6 di 8 penalità.\\ \tt Riprova! 
	  } 
	   } 
	   } 
	 & \\ 
 

    \theutterance \stepcounter{utterance}  

    & & & \multicolumn{4}{p{0.6\linewidth}}{\cellcolor[rgb]{1,0.85,0.72}{%
	\makecell[{{p{\linewidth}}}]{% 
	  \tt {\tiny [GM$\langle$B]}  
	 SPOSTA: P, (7, 1) 
	  } 
	   } 
	   } 
	  \\ 
 

    \theutterance \stepcounter{utterance}  

    & & & \multicolumn{2}{p{0.3\linewidth}}{\cellcolor[rgb]{0.95,0.95,0.95}{%
	\makecell[{{p{\linewidth}}}]{% 
	  \tt {\tiny [GM$|$GM]}  
	 SPOSTA: P, (7, 1) 
	  } 
	   } 
	   } 
	 & & \\ 
 

    \theutterance \stepcounter{utterance}  

    & & & \multicolumn{2}{p{0.3\linewidth}}{\cellcolor[rgb]{0.95,0.95,0.95}{%
	\makecell[{{p{\linewidth}}}]{% 
	  \tt {\tiny [GM$|$GM]}  
	 Spostato 'P' in (7,1) con successo. La tua griglia aggiornata appare così:\\ \tt ```\\ \tt  1234567\\ \tt ╔═╤═╤═╤═╗\\ \tt ║L│C│◌│P║ 1\\ \tt ║◌└─┼─┴─╢ 2\\ \tt ║◌◌◌│◌◌◌║ 3\\ \tt ╟──┬┴───╢ 4\\ \tt ║◌◌│◌◌◌◌║ 5\\ \tt ║◌◌│◌◌◌◌║ 6\\ \tt ║◌◌│◌◌◌◌║ 7\\ \tt ╚══╧════╝\\ \tt ```\\ \tt  
	  } 
	   } 
	   } 
	 & & \\ 
 

    \theutterance \stepcounter{utterance}  

    & & \multicolumn{4}{p{0.6\linewidth}}{\cellcolor[rgb]{0.9,0.9,0.9}{%
	\makecell[{{p{\linewidth}}}]{% 
	  \tt {\tiny [A$\langle$GM]}  
	 Il tuo messaggio è stato inoltrato all'altro giocatore.\\ \tt \\ \tt \\ \tt Stai attualmente giocando il round 8 di massimo 12.\\ \tt Avete accumulato collettivamente 6 di 8 penalità.\\ \tt L'altro giocatore ha spostato un oggetto sulla sua griglia.\\ \tt \\ \tt \\ \tt Qual è il tuo prossimo comando? 
	  } 
	   } 
	   } 
	 & \\ 
 

    \theutterance \stepcounter{utterance}  

    & \multicolumn{4}{p{0.6\linewidth}}{\cellcolor[rgb]{0.8,1,0.9}{%
	\makecell[{{p{\linewidth}}}]{% 
	  \tt {\tiny [A$\rangle$GM]}  
	 COMUNICA: Perfetto! I miei oggetti sono già nelle posizioni finali: L in (1,1), C in (3,1), e P in (5,1). Hai completato anche i tuoi spostamenti? Se sì, siamo pronti a finire! 
	  } 
	   } 
	   } 
	 & & \\ 
 

    \theutterance \stepcounter{utterance}  

    & & & \multicolumn{2}{p{0.3\linewidth}}{\cellcolor[rgb]{0.95,0.95,0.95}{%
	\makecell[{{p{\linewidth}}}]{% 
	  \tt {\tiny [GM$|$GM]}  
	 COMUNICA: Perfetto! I miei oggetti sono già nelle posizioni finali: L in (1,1), C in (3,1), e P in (5,1). Hai completato anche i tuoi spostamenti? Se sì, siamo pronti a finire! 
	  } 
	   } 
	   } 
	 & & \\ 
 

    \theutterance \stepcounter{utterance}  

    & & \multicolumn{4}{p{0.6\linewidth}}{\cellcolor[rgb]{0.9,0.9,0.9}{%
	\makecell[{{p{\linewidth}}}]{% 
	  \tt {\tiny [GM$\rangle$B]}  
	 Spostato 'P' in (7,1) con successo. La tua griglia aggiornata appare così:\\ \tt ```\\ \tt  1234567\\ \tt ╔═╤═╤═╤═╗\\ \tt ║L│C│◌│P║ 1\\ \tt ║◌└─┼─┴─╢ 2\\ \tt ║◌◌◌│◌◌◌║ 3\\ \tt ╟──┬┴───╢ 4\\ \tt ║◌◌│◌◌◌◌║ 5\\ \tt ║◌◌│◌◌◌◌║ 6\\ \tt ║◌◌│◌◌◌◌║ 7\\ \tt ╚══╧════╝\\ \tt ```\\ \tt Stai attualmente giocando il round 9 di massimo 12.\\ \tt Avete accumulato collettivamente 6 di 8 penalità.\\ \tt L'altro giocatore ha inviato questo messaggio:\\ \tt \\ \tt \\ \tt "Perfetto! I miei oggetti sono già nelle posizioni finali: L in (1,1), C in (3,1), e P in (5,1). Hai completato anche i tuoi spostamenti? Se sì, siamo pronti a finire!"\\ \tt \\ \tt \\ \tt Qual è il tuo prossimo comando? 
	  } 
	   } 
	   } 
	 & \\ 
 

    \theutterance \stepcounter{utterance}  

    & & & \multicolumn{4}{p{0.6\linewidth}}{\cellcolor[rgb]{1,0.85,0.72}{%
	\makecell[{{p{\linewidth}}}]{% 
	  \tt {\tiny [GM$\langle$B]}  
	 COMUNICA: Non ancora! Ho L in (1,1) correttamente, ma C è ancora in (2,1) e P è in (7,1). Devo spostare C da (2,1) a (3,1) per allinearmi con la tua configurazione. Lo faccio ora. 
	  } 
	   } 
	   } 
	  \\ 
 

    \theutterance \stepcounter{utterance}  

    & & & \multicolumn{2}{p{0.3\linewidth}}{\cellcolor[rgb]{0.95,0.95,0.95}{%
	\makecell[{{p{\linewidth}}}]{% 
	  \tt {\tiny [GM$|$GM]}  
	 COMUNICA: Non ancora! Ho L in (1,1) correttamente, ma C è ancora in (2,1) e P è in (7,1). Devo spostare C da (2,1) a (3,1) per allinearmi con la tua configurazione. Lo faccio ora. 
	  } 
	   } 
	   } 
	 & & \\ 
 

    \theutterance \stepcounter{utterance}  

    & & \multicolumn{4}{p{0.6\linewidth}}{\cellcolor[rgb]{0.9,0.9,0.9}{%
	\makecell[{{p{\linewidth}}}]{% 
	  \tt {\tiny [A$\langle$GM]}  
	 Il tuo messaggio è stato inoltrato all'altro giocatore.\\ \tt \\ \tt \\ \tt Stai attualmente giocando il round 9 di massimo 12.\\ \tt Avete accumulato collettivamente 6 di 8 penalità.\\ \tt L'altro giocatore ha inviato questo messaggio:\\ \tt \\ \tt \\ \tt "Non ancora! Ho L in (1,1) correttamente, ma C è ancora in (2,1) e P è in (7,1). Devo spostare C da (2,1) a (3,1) per allinearmi con la tua configurazione. Lo faccio ora."\\ \tt \\ \tt \\ \tt Qual è il tuo prossimo comando? 
	  } 
	   } 
	   } 
	 & \\ 
 

    \theutterance \stepcounter{utterance}  

    & \multicolumn{4}{p{0.6\linewidth}}{\cellcolor[rgb]{0.8,1,0.9}{%
	\makecell[{{p{\linewidth}}}]{% 
	  \tt {\tiny [A$\rangle$GM]}  
	 COMUNICA: Perfetto! Dopo che sposti anche P da (7,1) a (5,1), avremo tutti gli oggetti allineati. Ti aspetto per il movimento finale di P. 
	  } 
	   } 
	   } 
	 & & \\ 
 

    \theutterance \stepcounter{utterance}  

    & & & \multicolumn{2}{p{0.3\linewidth}}{\cellcolor[rgb]{0.95,0.95,0.95}{%
	\makecell[{{p{\linewidth}}}]{% 
	  \tt {\tiny [GM$|$GM]}  
	 COMUNICA: Perfetto! Dopo che sposti anche P da (7,1) a (5,1), avremo tutti gli oggetti allineati. Ti aspetto per il movimento finale di P. 
	  } 
	   } 
	   } 
	 & & \\ 
 

    \theutterance \stepcounter{utterance}  

    & & \multicolumn{4}{p{0.6\linewidth}}{\cellcolor[rgb]{0.9,0.9,0.9}{%
	\makecell[{{p{\linewidth}}}]{% 
	  \tt {\tiny [GM$\rangle$B]}  
	 Il tuo messaggio è stato inoltrato all'altro giocatore.\\ \tt \\ \tt \\ \tt Stai attualmente giocando il round 10 di massimo 12.\\ \tt Avete accumulato collettivamente 6 di 8 penalità.\\ \tt L'altro giocatore ha inviato questo messaggio:\\ \tt \\ \tt \\ \tt "Perfetto! Dopo che sposti anche P da (7,1) a (5,1), avremo tutti gli oggetti allineati. Ti aspetto per il movimento finale di P."\\ \tt \\ \tt \\ \tt Qual è il tuo prossimo comando? 
	  } 
	   } 
	   } 
	 & \\ 
 

    \theutterance \stepcounter{utterance}  

    & & & \multicolumn{4}{p{0.6\linewidth}}{\cellcolor[rgb]{1,0.85,0.72}{%
	\makecell[{{p{\linewidth}}}]{% 
	  \tt {\tiny [GM$\langle$B]}  
	 SPOSTA: C, (3, 1) 
	  } 
	   } 
	   } 
	  \\ 
 

    \theutterance \stepcounter{utterance}  

    & & & \multicolumn{2}{p{0.3\linewidth}}{\cellcolor[rgb]{0.95,0.95,0.95}{%
	\makecell[{{p{\linewidth}}}]{% 
	  \tt {\tiny [GM$|$GM]}  
	 SPOSTA: C, (3, 1) 
	  } 
	   } 
	   } 
	 & & \\ 
 

    \theutterance \stepcounter{utterance}  

    & & & \multicolumn{2}{p{0.3\linewidth}}{\cellcolor[rgb]{0.95,0.95,0.95}{%
	\makecell[{{p{\linewidth}}}]{% 
	  \tt {\tiny [GM$|$GM]}  
	 Penalità: (3,1) non è vuoto, ma contiene 'C'.\\ \tt \\ \tt Avete accumulato collettivamente 7 di 8 penalità.\\ \tt Riprova! 
	  } 
	   } 
	   } 
	 & & \\ 
 

    \theutterance \stepcounter{utterance}  

    & & \multicolumn{4}{p{0.6\linewidth}}{\cellcolor[rgb]{0.9,0.9,0.9}{%
	\makecell[{{p{\linewidth}}}]{% 
	  \tt {\tiny [GM$\rangle$B]}  
	 Penalità: (3,1) non è vuoto, ma contiene 'C'.\\ \tt \\ \tt Avete accumulato collettivamente 7 di 8 penalità.\\ \tt Riprova! 
	  } 
	   } 
	   } 
	 & \\ 
 

    \theutterance \stepcounter{utterance}  

    & & & \multicolumn{4}{p{0.6\linewidth}}{\cellcolor[rgb]{1,0.85,0.72}{%
	\makecell[{{p{\linewidth}}}]{% 
	  \tt {\tiny [GM$\langle$B]}  
	 COMUNICA: Scusa, c'è stata confusione. Guardando bene la mia griglia, ho L in (1,1), C in (2,1) e P in (7,1). Per allinearmi con te, devo spostare P da (7,1) a (5,1) e poi C da (2,1) a (3,1). Provo prima con P. 
	  } 
	   } 
	   } 
	  \\ 
 

    \theutterance \stepcounter{utterance}  

    & & & \multicolumn{2}{p{0.3\linewidth}}{\cellcolor[rgb]{0.95,0.95,0.95}{%
	\makecell[{{p{\linewidth}}}]{% 
	  \tt {\tiny [GM$|$GM]}  
	 COMUNICA: Scusa, c'è stata confusione. Guardando bene la mia griglia, ho L in (1,1), C in (2,1) e P in (7,1). Per allinearmi con te, devo spostare P da (7,1) a (5,1) e poi C da (2,1) a (3,1). Provo prima con P. 
	  } 
	   } 
	   } 
	 & & \\ 
 

    \theutterance \stepcounter{utterance}  

    & & \multicolumn{4}{p{0.6\linewidth}}{\cellcolor[rgb]{0.9,0.9,0.9}{%
	\makecell[{{p{\linewidth}}}]{% 
	  \tt {\tiny [A$\langle$GM]}  
	 Il tuo messaggio è stato inoltrato all'altro giocatore.\\ \tt \\ \tt \\ \tt Stai attualmente giocando il round 10 di massimo 12.\\ \tt Avete accumulato collettivamente 7 di 8 penalità.\\ \tt L'altro giocatore ha inviato questo messaggio:\\ \tt \\ \tt \\ \tt "Scusa, c'è stata confusione. Guardando bene la mia griglia, ho L in (1,1), C in (2,1) e P in (7,1). Per allinearmi con te, devo spostare P da (7,1) a (5,1) e poi C da (2,1) a (3,1). Provo prima con P."\\ \tt \\ \tt \\ \tt Qual è il tuo prossimo comando? 
	  } 
	   } 
	   } 
	 & \\ 
 

    \theutterance \stepcounter{utterance}  

    & \multicolumn{4}{p{0.6\linewidth}}{\cellcolor[rgb]{0.8,1,0.9}{%
	\makecell[{{p{\linewidth}}}]{% 
	  \tt {\tiny [A$\rangle$GM]}  
	 COMUNICA: Va bene, ti aspetto! Una volta che avrai spostato P in (5,1) e poi C in (3,1), avremo la stessa configurazione: L in (1,1), C in (3,1), P in (5,1). Dimmi quando hai finito! 
	  } 
	   } 
	   } 
	 & & \\ 
 

    \theutterance \stepcounter{utterance}  

    & & & \multicolumn{2}{p{0.3\linewidth}}{\cellcolor[rgb]{0.95,0.95,0.95}{%
	\makecell[{{p{\linewidth}}}]{% 
	  \tt {\tiny [GM$|$GM]}  
	 COMUNICA: Va bene, ti aspetto! Una volta che avrai spostato P in (5,1) e poi C in (3,1), avremo la stessa configurazione: L in (1,1), C in (3,1), P in (5,1). Dimmi quando hai finito! 
	  } 
	   } 
	   } 
	 & & \\ 
 

    \theutterance \stepcounter{utterance}  

    & & \multicolumn{4}{p{0.6\linewidth}}{\cellcolor[rgb]{0.9,0.9,0.9}{%
	\makecell[{{p{\linewidth}}}]{% 
	  \tt {\tiny [GM$\rangle$B]}  
	 Il tuo messaggio è stato inoltrato all'altro giocatore.\\ \tt \\ \tt \\ \tt Stai attualmente giocando il round 11 di massimo 12.\\ \tt Avete accumulato collettivamente 7 di 8 penalità.\\ \tt L'altro giocatore ha inviato questo messaggio:\\ \tt \\ \tt \\ \tt "Va bene, ti aspetto! Una volta che avrai spostato P in (5,1) e poi C in (3,1), avremo la stessa configurazione: L in (1,1), C in (3,1), P in (5,1). Dimmi quando hai finito!"\\ \tt \\ \tt \\ \tt Qual è il tuo prossimo comando? 
	  } 
	   } 
	   } 
	 & \\ 
 

    \theutterance \stepcounter{utterance}  

    & & & \multicolumn{4}{p{0.6\linewidth}}{\cellcolor[rgb]{1,0.85,0.72}{%
	\makecell[{{p{\linewidth}}}]{% 
	  \tt {\tiny [GM$\langle$B]}  
	 SPOSTA: P, (5, 1) 
	  } 
	   } 
	   } 
	  \\ 
 

    \theutterance \stepcounter{utterance}  

    & & & \multicolumn{2}{p{0.3\linewidth}}{\cellcolor[rgb]{0.95,0.95,0.95}{%
	\makecell[{{p{\linewidth}}}]{% 
	  \tt {\tiny [GM$|$GM]}  
	 SPOSTA: P, (5, 1) 
	  } 
	   } 
	   } 
	 & & \\ 
 

    \theutterance \stepcounter{utterance}  

    & & & \multicolumn{2}{p{0.3\linewidth}}{\cellcolor[rgb]{0.95,0.95,0.95}{%
	\makecell[{{p{\linewidth}}}]{% 
	  \tt {\tiny [GM$|$GM]}  
	 Spostato 'P' in (5,1) con successo. La tua griglia aggiornata appare così:\\ \tt ```\\ \tt  1234567\\ \tt ╔═╤═╤═╤═╗\\ \tt ║L│C│P│◌║ 1\\ \tt ║◌└─┼─┴─╢ 2\\ \tt ║◌◌◌│◌◌◌║ 3\\ \tt ╟──┬┴───╢ 4\\ \tt ║◌◌│◌◌◌◌║ 5\\ \tt ║◌◌│◌◌◌◌║ 6\\ \tt ║◌◌│◌◌◌◌║ 7\\ \tt ╚══╧════╝\\ \tt ```\\ \tt  
	  } 
	   } 
	   } 
	 & & \\ 
 

    \theutterance \stepcounter{utterance}  

    & & \multicolumn{4}{p{0.6\linewidth}}{\cellcolor[rgb]{0.9,0.9,0.9}{%
	\makecell[{{p{\linewidth}}}]{% 
	  \tt {\tiny [A$\langle$GM]}  
	 Il tuo messaggio è stato inoltrato all'altro giocatore.\\ \tt \\ \tt \\ \tt Stai attualmente giocando il round 11 di massimo 12.\\ \tt Avete accumulato collettivamente 7 di 8 penalità.\\ \tt L'altro giocatore ha spostato un oggetto sulla sua griglia.\\ \tt \\ \tt \\ \tt Qual è il tuo prossimo comando? 
	  } 
	   } 
	   } 
	 & \\ 
 

    \theutterance \stepcounter{utterance}  

    & \multicolumn{4}{p{0.6\linewidth}}{\cellcolor[rgb]{0.8,1,0.9}{%
	\makecell[{{p{\linewidth}}}]{% 
	  \tt {\tiny [A$\rangle$GM]}  
	 COMUNICA: Siamo all'ultimo round! Hai completato tutti i tuoi spostamenti? Se ora hai L in (1,1), C in (3,1), e P in (5,1) come me, finito? 
	  } 
	   } 
	   } 
	 & & \\ 
 

    \theutterance \stepcounter{utterance}  

    & & & \multicolumn{2}{p{0.3\linewidth}}{\cellcolor[rgb]{0.95,0.95,0.95}{%
	\makecell[{{p{\linewidth}}}]{% 
	  \tt {\tiny [GM$|$GM]}  
	 COMUNICA: Siamo all'ultimo round! Hai completato tutti i tuoi spostamenti? Se ora hai L in (1,1), C in (3,1), e P in (5,1) come me, finito? 
	  } 
	   } 
	   } 
	 & & \\ 
 

    \theutterance \stepcounter{utterance}  

    & & \multicolumn{4}{p{0.6\linewidth}}{\cellcolor[rgb]{0.9,0.9,0.9}{%
	\makecell[{{p{\linewidth}}}]{% 
	  \tt {\tiny [GM$\rangle$B]}  
	 Spostato 'P' in (5,1) con successo. La tua griglia aggiornata appare così:\\ \tt ```\\ \tt  1234567\\ \tt ╔═╤═╤═╤═╗\\ \tt ║L│C│P│◌║ 1\\ \tt ║◌└─┼─┴─╢ 2\\ \tt ║◌◌◌│◌◌◌║ 3\\ \tt ╟──┬┴───╢ 4\\ \tt ║◌◌│◌◌◌◌║ 5\\ \tt ║◌◌│◌◌◌◌║ 6\\ \tt ║◌◌│◌◌◌◌║ 7\\ \tt ╚══╧════╝\\ \tt ```\\ \tt Stai attualmente giocando il round 12 di massimo 12.\\ \tt Avete accumulato collettivamente 7 di 8 penalità.\\ \tt L'altro giocatore ha inviato questo messaggio:\\ \tt \\ \tt \\ \tt "Siamo all'ultimo round! Hai completato tutti i tuoi spostamenti? Se ora hai L in (1,1), C in (3,1), e P in (5,1) come me, finito?"\\ \tt \\ \tt \\ \tt Qual è il tuo prossimo comando? 
	  } 
	   } 
	   } 
	 & \\ 
 

    \theutterance \stepcounter{utterance}  

    & & & \multicolumn{4}{p{0.6\linewidth}}{\cellcolor[rgb]{1,0.85,0.72}{%
	\makecell[{{p{\linewidth}}}]{% 
	  \tt {\tiny [GM$\langle$B]}  
	 COMUNICA: Non siamo ancora allineati! Ho L in (1,1), C in (2,1), e P in (3,1). Tu hai L in (1,1), C in (3,1), P in (5,1). Non riesco a completare tutti gli spostamenti in un solo comando finale. Purtroppo non credo riusciremo a finire perfettamente allineati! 
	  } 
	   } 
	   } 
	  \\ 
 

    \theutterance \stepcounter{utterance}  

    & & & \multicolumn{2}{p{0.3\linewidth}}{\cellcolor[rgb]{0.95,0.95,0.95}{%
	\makecell[{{p{\linewidth}}}]{% 
	  \tt {\tiny [GM$|$GM]}  
	 COMUNICA: Non siamo ancora allineati! Ho L in (1,1), C in (2,1), e P in (3,1). Tu hai L in (1,1), C in (3,1), P in (5,1). Non riesco a completare tutti gli spostamenti in un solo comando finale. Purtroppo non credo riusciremo a finire perfettamente allineati! 
	  } 
	   } 
	   } 
	 & & \\ 
 

    \theutterance \stepcounter{utterance}  

    & & & \multicolumn{2}{p{0.3\linewidth}}{\cellcolor[rgb]{0.95,0.95,0.95}{%
	\makecell[{{p{\linewidth}}}]{% 
	  \tt {\tiny [GM$|$GM]}  
	 Maximum number of rounds reached 
	  } 
	   } 
	   } 
	 & & \\ 
 

    \theutterance \stepcounter{utterance}  

    & & & \multicolumn{2}{p{0.3\linewidth}}{\cellcolor[rgb]{0.95,0.95,0.95}{%
	\makecell[{{p{\linewidth}}}]{% 
	  \tt {\tiny [GM$|$GM]}  
	 * success: True\\ \tt * lose: False\\ \tt * aborted: False\\ \tt -------\\ \tt * Shifts: 2.00\\ \tt * Max Shifts: 4.00\\ \tt * Min Shifts: 2.00\\ \tt * End Distance Sum: 0.00\\ \tt * Init Distance Sum: 13.65\\ \tt * Expected Distance Sum: 12.57\\ \tt * Penalties: 7.00\\ \tt * Max Penalties: 8.00\\ \tt * Rounds: 12.00\\ \tt * Max Rounds: 12.00\\ \tt * Object Count: 3.00\\ \tt  
	  } 
	   } 
	   } 
	 & & \\ 
 

    \theutterance \stepcounter{utterance}  

    & & & \multicolumn{2}{p{0.3\linewidth}}{\cellcolor[rgb]{0.95,0.95,0.95}{%
	\makecell[{{p{\linewidth}}}]{% 
	  \tt {\tiny [GM$|$GM]}  
	 * Main Score: 28.57\\ \tt \\ \tt -------\\ \tt * Distance Score: 1.00\\ \tt * Consistency Score: 1.00\\ \tt * Coverage Score: 1.00\\ \tt * Penalty Score: 0.29\\ \tt * Alternative Penalty Score: 0.44\\ \tt * Alternative Main Score: 44.44\\ \tt \\ \tt -------\\ \tt * Shifts: 2.00\\ \tt * Max Shifts: 4.00\\ \tt * Min Shifts: 2.00\\ \tt * End Distance Sum: 0.00\\ \tt * Init Distance Sum: 13.65\\ \tt * Expected Distance Sum: 12.57\\ \tt * Penalties: 7.00\\ \tt * Max Penalties: 8.00\\ \tt * Rounds: 12.00\\ \tt * Max Rounds: 12.00\\ \tt * Object Count: 3.00\\ \tt  
	  } 
	   } 
	   } 
	 & & \\ 
 

\end{supertabular}
}

\end{document}
